\section{Glosario de s\'imbolos}
\label{sec:glosario}

Para facilitar la lectura y unificar la notaci\'on utilizada a lo largo del
documento, en el Cuadro~\ref{tab:glosario-simbolos} se resume la nomenclatura
formal de par\'ametros y variables.

\begin{table}[!ht]
\centering
\small
\begin{tabularx}{\textwidth}{@{}l X@{}}
\toprule
S\'imbolo & Descripci\'on formal \\
\midrule
$G, g, q$ & Grupo algebraico c\'iclico (\texttt{ristretto255}), generador can\'onico $g$ y orden primo $q = 2^{252}+\dots$ \\
$N, i, j$ & Cantidad total de registros ($500 \dots 5000$); \'indice solicitado $i$; e \'indice de iteraci\'on $j$. \\
$K_j, C_j$ & Clave sim\'etrica AES-256 del registro $j$ y texto cifrado de 256\,B m\'as tag ($272$\,B en total). \\
$\beta, v$ & Escalar secreto ef\'imero del emisor ($\beta \in \mathbb{Z}_q$) y par\'ametro p\'ublico de sesi\'on ($v = g^\beta \in G$). \\
$\alpha, u$ & Escalar secreto del receptor ($\alpha \in \mathbb{Z}_q$) y elemento de consulta enmascarado ($u = g^\alpha v^{-i} \in G$). \\
$sid$ & Identificador fresco de sesi\'on de 128 bits (16\,bytes, representado en hexadecimal). \\
$\mathsf{ctx}$ & Cadena de contexto global inmutable con prefijos de longitud: cat\'alogo, versi\'on y grupo. \\
$AAD_{\mathrm{cat}}, AAD_{\mathrm{ot}}$ & Datos asociados: $\mathsf{ctx}\Vert j$ para registros y $\mathsf{ctx}\Vert sid\Vert j$ para encapsulaciones. \\
$u_j, Z_j, k_j, E_j$ & Clave p\'ublica impl\'icita ($u\cdot v^j$), secreto DH ($u_j^\beta$), clave derivada HKDF y clave envuelta. \\
$\mathrm{HKDF}, \mathrm{AEAD}$ & Funci\'on de derivaci\'on HKDF-SHA256 (RFC~5869) y cifrador autenticado AES-256-GCM (SP~800-38D). \\
\bottomrule
\end{tabularx}
\caption{Glosario formal de s\'imbolos y variables del protocolo.}
\label{tab:glosario-simbolos}
\end{table}
